% Verbatim copy of /Users/alingalatan/workspaces/grothe/handoff/manuscript/legacy/ch25-parts/87-phase-lift-screening.tex (file dated 2026-10-25). Only this line was added.
% Consolidated source module (GPT phase-cut experiment, 22 October 2026). Owner: ch25-weak-k.tex
\subsection{A screening theorem and a finite obstruction for phase cuts}
\label{sec:phase-lift-screening-20261022}

\paragraph{Status and scope.}
The results below are \st{P,C} restrictions on two particular uses of
phase-correlation constraints. They give no new lower bound for $\KG$.
They neither settle the actual-response inequality nor place a ceiling
on all Gaussian configurations or on the complete certificate cover.

Let
\[
 H=\begin{pmatrix}0&-i&i&1\\i&0&-i&1\\-i&i&0&1\\1&1&1&0\end{pmatrix}.
\]
The phase inequality $u^*Hu\le6$ is due to Jarre--Lieder--Liu--Lu;
the small Gram certificate below appears in Sinjorgo--Sotirov--Anjos,
\emph{Cuts and semidefinite liftings for the complex cut polytope},
Section~5, Lemma~12 of the Optimization Online preprint
(arXiv:2402.04731).
For $u=(1,z_1,z_2,z_3)$ and
\[
 v=(1,z_1,z_2,z_3,\overline z_1z_2,\overline z_1z_3)^T,
 \quad a=(2,0,i,-1,-i,-1)^T,\quad b=(0,0,i,-i,-1,1)^T,
\]
one has $6-u^*Hu=(|a^*v|^2+|b^*v|^2)/2$ on the torus, and $H^2=3I$.

\begin{proposition}[Retained-energy screening, \st{P,C}]
Let $G$ be any complex correlation matrix of order four, and assign
either response label to each index. Retain any collection of chaos
blocks of the two responses, with retained energies $e_+,e_-$.
If $M$ is their transferred covariance matrix, then
\[
 \operatorname{tr}(HM)\le\sqrt3\sum_{j=1}^4e_{\sigma_j}.
\]
Thus $e_+,e_-\le\sqrt3/2$ makes the retained phase inequality automatic,
uniformly in the Gaussian dimension and configuration. The statement
also holds after diagonal unitary conjugation, permutation, or complex
conjugation of $H$.
\end{proposition}
\begin{proof}
Each retained block is a Schur product of two positive semidefinite
matrices: the Gaussian chaos kernel and the label-selected response
Gram matrix. Hence $M\succeq0$, with the indicated diagonal energies.
Since $H\preceq\sqrt3I$, the trace inequality follows. The transforms
of $H$ preserve both its spectrum and its phase bound.
\end{proof}

For $e_+=e_-=1707/2000$, the margin is strictly greater than
$2421/28000$, using $\sqrt3<97/56$. This also rules out separation by
the retained inequality with a nonnegative omitted-tail allowance.
It does not rule out a signed treatment of resolved tail blocks.
The threshold is sharp for arbitrary positive semidefinite matrices
with the stated diagonal bound, as shown by $e(I+H/\sqrt3)$; this
sharpness assertion does not concern actual bounded responses.

\begin{proposition}[One exact twelve-variable lift obstruction, \st{P,C}]
There is a formal charge-one response profile with blocks
\[
 A_0=\frac{289}{400}\begin{pmatrix}1&1\\1&1\end{pmatrix},\qquad
 A_1=\begin{pmatrix}131/1000&11/1000\\11/1000&131/1000\end{pmatrix},\qquad
 A_{150}=\frac{293}{2000}I_2,
\]
and all other blocks zero, satisfying positivity, Parseval, the
first-chaos bound and the cubic-radius bound. At the specific
twelve-variable configuration below it has a feasible square-free
degree-two moment lift, with eigenvalue margin greater than $9/1000$.
Its formal objective is $0.714305$ for $\epsilon=149/200$, and
$0.7141796$ for $\epsilon=1891/2500$.
\end{proposition}
\begin{proof}
Use the six rational unit vectors
\[
 u(b,c)=\left(\frac{1-b^2-c^2}{1+b^2+c^2},
              \frac{2b+2ic}{1+b^2+c^2}\right),\qquad
 (b,c)\in\{(0,0),(1/2,0),(0,1/2),(1/2,1/2),(2,0),(0,2)\}.
\]
Duplicate them and assign labels $+$ and $-$ to the respective copies.
Their Gram matrix has squared off-site correlations at most $9/10$.
The lift uses
$\mathcal B=\{0,e_1,\ldots,e_{11}\}\cup\{e_i+e_j:i<j\}$,
of cardinality $67$.

The supplied rational certificate constructs a moment matrix $M^0$
for the first two blocks, with diagonal one, and a rational matrix $L$.
For $R=M^0-LL^*$, every row satisfies
\[
 R_{ii}-\sum_{j\ne i}(|\Re R_{ij}|+|\Im R_{ij}|)>1/100.
\]
Thus $M^0\succ(1/100)I$. Adding the last block changes each prescribed
moment by at most $(293/2000)(9/10)^{150}$; at duplicate sites the
cross-label change is zero. Copying these changes through the lift
has operator norm at most $67(293/2000)(9/10)^{150}<1/1000$.
This proves the asserted margin. All comparisons and the two objective
values are checked over the rationals by \texttt{verify\_all.py}.
\end{proof}

\paragraph{Interpretation.}
The formal profile is not a bounded response pair: its finite nonconstant
Gaussian chaos expansion would be an unbounded polynomial. Passing this
lift does not assert membership in the phase body. The result only shows
that the selected configuration and moment basis cannot distinguish this
infeasible profile. Its compatibility with the full manuscript relaxation
has not been established. Floating searches at other configurations remain
\st{F}; no universal conclusion is drawn from them. No conjectural
exactness of a moment lift, fresh-buffer independence, or conditional
Gaussianity assumption is used.


\paragraph{Provenance and IDE replay (22 October).} Package \code{grothendieck-asymmetric-phase-tests.zip}, archived with the replay log in
\pathref{gpt-reviews/phase-tests-20261022/}; its standard-library verifier PASSES under \code{python3 -I} from an unrelated directory (0.5\,s): the four-phase
SOS identity and $H^2=3I$, the rational PSD profile, the two objective comparisons, the $67\times67$ rational completion with margin $1/100$, the
degree-301 tail perturbation $<1/1000$, and the manifest. IDE check: $H$ Hermitian with $H^2=3I$ has spectrum $\pm\sqrt3$, so $H\preceq\sqrt3I$ and the
screening bound $\operatorname{tr}(HM)\le\sqrt3\operatorname{tr}M$ is immediate once $M\succeq0$ (Schur products of PSD blocks). Both results are negative
(stop rules for the phase-cut route); no bound changes.
